\documentclass[11pt]{article}
\usepackage[margin=1in]{geometry}
\usepackage{amsmath,amssymb}
\usepackage{listings}
\usepackage[T1]{fontenc}
\usepackage[utf8]{inputenc}
\usepackage{ctex}

\title{Adam 与 AdamW：原理与实现笔记}
\author{}
\date{}

\begin{document}
\maketitle

\section{Adam 在做什么}
Adam 结合了两种思想：
\begin{itemize}
  \item 动量（梯度一阶矩），
  \item 自适应步长缩放（梯度二阶矩）。
\end{itemize}

在第 $t$ 步，记梯度为 $g_t = \nabla_\theta L_t$。

\section{Adam 更新公式}
\begin{align}
  m_t &= \beta_1 m_{t-1} + (1-\beta_1) g_t, \\
  v_t &= \beta_2 v_{t-1} + (1-\beta_2) g_t^2.
\end{align}

偏置修正：
\begin{align}
  \hat m_t &= \frac{m_t}{1-\beta_1^t}, \\
  \hat v_t &= \frac{v_t}{1-\beta_2^t}.
\end{align}

参数更新：
\begin{equation}
  \theta_{t+1} = \theta_t - \alpha \frac{\hat m_t}{\sqrt{\hat v_t} + \varepsilon}.
\end{equation}

\section{AdamW（解耦权重衰减）}
AdamW 在 Adam 主更新之后，\emph{额外}执行权重衰减：
\begin{equation}
  \theta \leftarrow \theta - \alpha\lambda\theta.
\end{equation}

这与把 $\lambda\theta$ 直接加到梯度中的做法不同。

\section{实现清单}
对每个参数张量：
\begin{enumerate}
  \item 首次初始化状态：\texttt{step}、\texttt{exp\_avg}、\texttt{exp\_avg\_sq}。
  \item 更新一阶矩与二阶矩。
  \item 若 \texttt{correct\_bias=True}，执行偏置修正。
  \item 按 Adam 公式做主更新。
  \item 再做 AdamW 的解耦权重衰减。
\end{enumerate}

\section{PyTorch 风格核心代码}
\begin{lstlisting}[language=Python,basicstyle=\ttfamily\small]
beta1, beta2 = group["betas"]
eps = group["eps"]
lr = group["lr"]
wd = group["weight_decay"]

state = self.state[p]
if len(state) == 0:
    state["step"] = 0
    state["exp_avg"] = torch.zeros_like(p.data)
    state["exp_avg_sq"] = torch.zeros_like(p.data)

exp_avg = state["exp_avg"]
exp_avg_sq = state["exp_avg_sq"]
state["step"] += 1
t = state["step"]

exp_avg.mul_(beta1).add_(grad, alpha=1 - beta1)
exp_avg_sq.mul_(beta2).addcmul_(grad, grad, value=1 - beta2)

denom = exp_avg_sq.sqrt().add_(eps)
step_size = lr
if group["correct_bias"]:
    bc1 = 1 - beta1 ** t
    bc2 = 1 - beta2 ** t
    step_size = lr * (bc2 ** 0.5) / bc1

p.data.addcdiv_(exp_avg, denom, value=-step_size)

if wd > 0:
    p.data.add_(p.data, alpha=-lr * wd)
\end{lstlisting}

\section{常见坑}
\begin{itemize}
  \item 把 L2 正则混进梯度，而不是使用解耦衰减。
  \item 忘记递增 \texttt{step}。
  \item 偏置修正公式写错。
  \item \texttt{return} 后残留不可达代码。
\end{itemize}

\end{document}
